% ************************** Thesis Abstract *****************************
% Use `abstract' as an option in the document class to print only the titlepage and the abstract.
\begin{abstract}
Nanoparticle (NP)-based drug delivery systems offer a promising approach for treating cancers with poor prognosis, such as pancreatic cancer. However, nanocarrier stability and cellular uptake are strongly influenced by haemodynamic shear stresses encountered during blood circulation, effects that conventional static in vitro assays fail to capture. There is a need for compact, automated platforms capable of replicating physiological flow conditions for preclinical NP characterisation.

This work presents the design, implementation, and characterisation of an ESP32-based closed-loop microfluidic recirculation system for biomimetic flow studies. The platform integrates a Bartels MP6 piezoelectric micropump, a Sensirion SLF3S thermal mass flow sensor, and a 12-bit DAC for amplitude control, all coordinated by a FreeRTOS dual-task firmware architecture. A discrete PID controller operating at \SI{10}{\hertz} regulates the flow rate by dynamically adjusting the pump drive amplitude, with anti-windup protection and real-time safety monitoring. A Python PyQt5 graphical user interface provides live flow visualisation, wall shear rate computation based on user-defined tube parameters, and CSV data export.

System characterisation demonstrated that the platform achieves stable, programmable flow rates covering the physiological shear rate range from venous to arteriolar conditions. The closed-loop controller maintains the target flow rate with low steady-state error over sustained operation periods exceeding \SI{30}{\minute}, with effective bubble management via an integrated bubble trap. The platform represents a significant improvement over the previous manual prototype, offering automated control, real-time monitoring, and reproducible experimental conditions at a fraction of the cost of commercial syringe pump systems. Future work will integrate the system with endothelial cell cultures for NP adhesion and uptake studies under controlled shear.
\end{abstract}
